
%#############################################################
%MAIN file for  project
%#############################################################
\documentclass[12pt]{report}
%############################################################
%###### Package 'Adding' 
%###### This package contains some tools to generate index, glossary
%###### nomenclature, acronym
%###### Author  : Nicolas LE GUERROUE
%###### Contact : nicolasleguerroue@gmail.com
%############################################################
%\ProvidesPackage{Utils}[2013/01/13 Utils Package]
%############################################################
%##### Glossaries
\RequirePackage[xindy, acronym, nomain, toc]{glossaries}
\RequirePackage{ifthen} 
\RequirePackage{nomencl}  %nomenclature
\RequirePackage{lmodern}
\RequirePackage[T1]{fontenc}
\RequirePackage[utf8]{inputenc}     %UTF-8 encoding
\RequirePackage{graphicx}           %Images
\RequirePackage{caption}            %légende
\RequirePackage{textcomp}           %special characters

\RequirePackage{lastpage}           %page counter
%\RequirePackage{fullpage}          %entire page
\RequirePackage{float}              %image floating
\RequirePackage{wrapfig}
\RequirePackage{subcaption}         %Subcaption
\RequirePackage{listings}           %new env such as Bash, Python
%\RequirePackage{tcolorbox}          %messageBox
\RequirePackage[many]{tcolorbox}
\RequirePackage{color}              %colors
\RequirePackage{geometry}
\RequirePackage{listings}           %new env such as Bash, Python
\RequirePackage{xcolor}             %define new colors
\RequirePackage{xparse}
\RequirePackage{amssymb,amsthm}     %math
\RequirePackage{multirow}
\RequirePackage{tabularx}           %use for csv file [table]

\RequirePackage{bookmark}           %use for hyperref
\RequirePackage{makeidx}            %make index

\makeindex
\makeglossaries
\makenomenclature

%##### Convention [Nomenclature]
  \renewcommand{\nomgroup}[1]{%
  \item[\bfseries
  \ifthenelse{\equal{#1}{P}}{Constantes physiques}{%
  \ifthenelse{\equal{#1}{O}}{Autres symboles}{%
  \ifthenelse{\equal{#1}{N}}{Nombres spéciaux}{ 
  \ifthenelse{\equal{#1}{A}}{Amplificateurs Opérationnels}{ 
  \ifthenelse{\equal{#1}{M}}{Mécanique}{ 
  \ifthenelse{\equal{#1}{E}}{Électronique}{}}}}}}%
  ]}

  \renewcommand{\glossary}[1]{%
  \gls{#1}
}
% Add unit on convention [nomenclature]
%----------------------------------------------
\newcommand{\addUnit}[1]{%
\renewcommand{\nomentryend}{\hspace*{\fill}#1}}
%----------------------------------------------%############################################################
%###### Package 'name' 
%###### This package contains ...
%###### Author  : Nicolas LE GUERROUE
%###### Contact : nicolasleguerroue@gmail.com
%############################################################
\RequirePackage{lmodern}
\RequirePackage[T1]{fontenc}
\RequirePackage[utf8]{inputenc}     %UTF-8 encoding
\RequirePackage[many]{tcolorbox}
%...

%1 : txt color
%2 : bg color 
%3 : text
\newtcbox{\badge}[3][red]{
  on line, 
  arc=2pt,
  colback=#3,
  colframe=#3,
  fontupper=\color{#2},
  boxrule=1pt, 
  boxsep=0pt,
  left=6pt,
  right=6pt,
  top=2pt,
  bottom=2pt
}

%################################################################%############################################################
%###### Package 'Chapter_1' 
%###### This package contains first model of chapter
%###### Author  : Nicolas LE GUERROUE
%###### Contact : nicolasleguerroue@gmail.com
%############################################################
% \RequirePackage[explicit]{titlesec} 


% \titleformat{\chapter}
%   {\gdef\chapterlabel{}
%    \normalfont\sffamily\Huge\bfseries\scshape}
%   {\gdef\chapterlabel{\thechapter)\ }}{0pt}
%   {\begin{tikzpicture}[remember picture,overlay]
%     \node[yshift=-2cm] at (current page.north west)
%       {\begin{tikzpicture}[remember picture, overlay]
%         \draw[fill=MediumBlue] (0,0) rectangle
%           (\paperwidth,2cm);
%         \node[anchor=east,xshift=.9\paperwidth,rectangle,
%               rounded corners=20pt,inner sep=11pt,
%               fill=white]
%               {\color{black}#1};%\chapterlabel
%           \node[anchor=west,yshift=1cm,xshift=2cm,inner sep=11pt,
%               fill=MediumBlue]
%               {\color{white}{\large Synthèse Latex}};%\chapterlabel
%        \end{tikzpicture}
%       };
%    \end{tikzpicture}
%   }
%   \titlespacing*{\chapter}{0pt}{50pt}{-60pt}%############################################################
%###### Package 'Color' 
%###### This package contains some colors
%###### Author  : Nicolas LE GUERROUE
%###### Contact : nicolasleguerroue@gmail.com
%############################################################
\RequirePackage{color}              %colors
%################################################################
\definecolor{pathColor}{rgb}{0.858, 0.188, 0.478}
\definecolor{fileColor}{rgb}{0.858, 0.188, 0.478}
\definecolor{water}{rgb}{0.858, 0.188, 0.478}
\definecolor{LightBlue}{RGB}{66, 163, 251}
\definecolor{DarkBlue}{RGB}{36, 100, 176}
\definecolor{LightGray}{gray}{.94}
\definecolor{DarkGray}{gray}{.172}
\definecolor{Orange}{RGB}{229, 133, 3}
\definecolor{MediumBlue}{RGB}{38, 119, 193}


\newcommand{\colors}[2]{
{\color{#1}{#2}}
}


%------------------Section Default Setting-------------%############################################################
%###### Package 'Electronic' 
%###### This package contains some tools to generate electornic circuits
%###### Author  : Nicolas LE GUERROUE
%###### Contact : nicolasleguerroue@gmail.com
%############################################################

\RequirePackage{tikz-timing}

\RequirePackage{graphics} %include figures
\RequirePackage{graphicx} %include figures
\RequirePackage{pgf,tikz}
\RequirePackage{circuitikz}
\usetikzlibrary{babel}  %allow to use tikz library with babel

\RequirePackage{ifthen}    %use if then 
 

%############################ Settings ##############################
\tikzset{%
    timing/table/axis/.style={->,>=latex},
    timing/table/axis ticks/.style={},   
}

%Direction of some device such as resistor, led...
%1.5 is the minimum length of device according my runs
\newcommand{\Up}{0,1.5}
\newcommand{\Down}{0,-1.5}
\newcommand{\Right}{1.5,0}
\newcommand{\Left}{-1.5,0}


%###### Length of components
\newcommand{\bipolesLength}[1]{#1cm}     %begin default size
%Length update
\newcommand{\setBipolesLength}[1]{
    \renewcommand{\bipolesLength}{#1}
    \ctikzset{bipoles/length=\bipolesLength cm}
}%End \setBipolesLength

%############ Mirrors and inverting
\newcommand{\Mirror}{}
\newcommand{\Invert}{}

%Update
\newcommand{\setMirror}[1]{
    \renewcommand{\Mirror}{,mirror}
}%End \setMirror

\newcommand{\setNoMirror}[1]{
    \renewcommand{\Mirror}{}
}%End \setMirror

\newcommand{\setInvert}[1]{
    \renewcommand{\Invert}{,Invert}
}%End \setInvert

\newcommand{\setNoInvert}[1]{
    \renewcommand{\Invert}{}
}%End \setInvert


%############## Rotate ###########
\newcommand{\rotate}{0}

%Update
\newcommand{\setRotate}[1]{
    \renewcommand{\rotate}{#1}
}%End \setRotate


%####################### Colors
%default colors of border colors and background colors
\newcommand{\deviceBorderColor}{black}
\newcommand{\deviceBackgroundColor}{white}

%Update
\newcommand{\setDeviceBorderColor}[1]{
    \renewcommand{\deviceBorderColor}{#1}
    \renewcommand{\deviceBackgroundColor}{white} %reset bg color
}%End \setDeviceBorderColor

\newcommand{\setDeviceBackgroundColor}[1]{
    \renewcommand{\deviceBorderColor}{black} %reset bg color
    \renewcommand{\deviceBackgroundColor}{#1}
}%\setDeviceBackgroundColor

%Reset
\newcommand{\resetColors}{
    \renewcommand{\deviceBorderColor}{black} %reset border color
    \renewcommand{\deviceBackgroundColor}{white} %reset bg color
}
%####################################################################
%############## draw device #########################################

%rotation
%color
%anchor

%\ifthenelse{\equal{#1}{0}}{A.}{no A.}
%Init

\begin{comment}
    @begin
    @command \addLogicGate
    @des 
    Cette commande permet de dessiner une porte logique à double entrée. Pour dessinder une porte inverseuse, utiliser la commande \addNotGate
    @sed
    @input Coordonnées de la porte en (x,y) sans parenthèse
    @input Référence de la porte pour s'accrocher aux entrées et sorties
    @input Type de la porte [nand, nor, or, and, or, xor]
    @input Label de sortie (laisser vide si absence de label souhaité)
    @input Label de l'entrée 1 (laisser vide si absence de label souhaité)
    @input Label de l'entrée 2 (laisser vide si absence de label souhaité)
    @input Nom de la porte [NOR1, AND1...]
    @begin_example 
    \addLogicGate{5,5}{logicgate}{nand}{S}{A}{B}{L1}
    @end_example
    @end
    \end{comment}

\newcommand{\addLogicGate}[7] {
    %\Colors
    \ifthenelse{\equal{\deviceBorderColor}  {black}}
    {\draw (#1)         node (#2) [rotate=\rotate,xshift=0cm,fill=\deviceBackgroundColor,#3 port] {#7}}%if equal to black
    {\draw (#1)         node (#2) [rotate=\rotate,xshift=0cm,color=\deviceBorderColor,#3 port] {#7}}

    (#2.out)  node      [anchor=south west, yshift=-0.3cm] {#4}
    (#2.in 1) node (A1)     [anchor=east,xshift=0cm,yshift=+0.3cm]       {#5}
    (#2.in 2) node (B1)     [anchor=east,xshift=0cm,yshift=+0.3cm]       {#6};
}


\newenvironment{schema}[1]
{
    \begin{center}
        \makeatletter
        \def\@captype{figure}
        \makeatother
        \newcommand{\TitleSchema}{#1}%use var to print title 

        
        %\shorthandoff{:;!?} %Compulsory if frenchb package is used (from babel)
            
        \begin{tikzpicture}
            %\setGraphic %command to display with frenchb babel
    }
    { 
        \end{tikzpicture}
   % \caption{\TitleSchema}
    \end{center}
}


\newenvironment{numeric}[1]
{
\begin{center}
    \makeatletter
    \def\@captype{figure}
    \makeatother
    \newcommand{\TitleNumeric}{#1}%use var to print title 
\begin{tikztimingtable}
}
{
\end{tikztimingtable}%
\caption{\TitleNumeric}
\end{center}
}


%cood
%name device
%type (npn, pnp)
%B
%C
%E
\newcommand{\addTransistor}[6] {

    \ifthenelse{\equal{\deviceBorderColor}  {black}}
    {\draw (#1)         node (#2) [xshift=0cm,fill=\deviceBackgroundColor,#3] {}}%if equal to black
    {\draw (#1)         node (#2) [xshift=0cm,color=\deviceBorderColor,#3] {}}

    (#2.B)  node      [anchor=south west, xshift=0cm, yshift=0cm] {#4} 
    (#2.C) node (A1)     [anchor=north,xshift=0.3cm,yshift=+0.1cm]       {#5}
    (#2.E) node (B1)     [anchor=south,xshift=0.3cm,yshift=0.1cm]       {#6};
}
%node 1
%node 2
%wire type
\newcommand{\addWire}[3] {
    \draw (#1) #3 (#2);
}%end addWire

\newcommand{\orthogonalWireA}{-|}
\newcommand{\orthogonalWireB}{|-}
\newcommand{\directWire}{--}

%coord
%label
%value
\newcommand{\addNode}[3] {
    \node (#2) at (#1) {#3};
}%end addWire


%posiition (x,y)
%color
%width in pt
\newcommand{\addPoint}[3] {
    \filldraw [#2] (#1) circle (#3pt);
}%end addWire

%posiition (x,y)
%name
%value
\newcommand{\addPower}[3] {
    \draw (#1) node (#2) [vcc] {#3};
}%end addPower

%posiition (x,y)
%name
%value
\newcommand{\addGround}[3] {
    \draw (#1) node (#2) [ground] {#3};
}%end addGround

%posiition depart(x,y)
%orientation de départ
%position arrivée
\newcommand{\addResistor}[4] {
    \draw (#1) to[R,l=$R$] +(#2) #4 (#3);
}%end addResistor

%posiition depart(x,y)
%orientation de départ
%position arrivée
%type de liaison
%nom de la led
\newcommand{\addLed}[5] {
    \ifthenelse{\equal{\deviceBorderColor}  {black}}
    {\draw (#1) to[leD,l_=#5,fill=\deviceBackgroundColor] +(#2) #4 (#3);}
    {\draw (#1) to[leD,l_=#5,color=\deviceBorderColor] +(#2) #4 (#3);}

    % {\draw (#1)         node (#2) [xshift=0cm,fill=\deviceBackgroundColor,#3] {}}%if equal to black
    % {\draw (#1)         node (#2) [xshift=0cm,color=\deviceBorderColor,#3] {}}

    %\draw (#1) to[leD,l_=#5] +(#2) #4 (#3);
}%end addResistor%############################################################
%###### Package 'Font' 
%###### This package contains some tools to set fonts
%###### Author  : Nicolas LE GUERROUE
%###### Contact : nicolasleguerroue@gmail.com
%############################################################
\RequirePackage{fontawesome}

\newcommand{\arobase}{\faAt}

   %############################################################
%###### Package 'Graphics' 
%###### This package contains some tools to create graphics 2D or 3D
%###### Author  : Nicolas LE GUERROUE
%###### Contact : nicolasleguerroue@gmail.com
%############################################################
%\ProvidesPackage{Utils}[2013/01/13 Utils Package]
%############################################################
\RequirePackage[T1]{fontenc}
\RequirePackage[utf8]{inputenc}     %UTF-8 encoding
\RequirePackage{csvsimple} 
\RequirePackage{tikz,pgfplots,pgf}  
\RequirePackage{version}            %use commented code

\RequirePackage{graphics} %include figures
\RequirePackage{graphicx} %include figures
\RequirePackage{caption}
\RequirePackage{subcaption} %Add
\RequirePackage{version}            %use commented code
\pgfplotsset{compat=1.7}
%###### Checking if babel is loaded
\makeatletter
\@ifpackageloaded{babel}
{% if the package was loaded
\newcommand{\setGraphic}{\shorthandoff{:;!?}} %Compulsory if frenchb package is used (from babel)
\frenchbsetup{StandardLists=true} %to include if using \RequirePackage[french]{babel} -> rounded list
}
{%else:
\newcommand{\setGraphic}{} %Compulsory if frenchb package is used (from babel)
}
\makeatother
%############################################################
%### WARNING : USE \shorthandoff{:;!?} before \begin{tikzpicture} 
%### environment
%############################################################
\begin{comment}
@begin
@env graphics
@des 
Cet environnement permet de tracer des courbes en 2D à partir de points, d'équations ou de fichier CSV
@sed
@input Largeur du graphique ]0;1] 
@input Hauteur du graphique ]0;1+] 
@input Valeur minimale en abscisse 
@input Valeur maximale en abscisse
@input Valeur minimale en ordonnée 
@input Valeur maximale en ordonnée
@input Légende de l'axe des abscisse
@input Légende de l'axe des ordonnées
@input Titre du graphique
@begin_example 
\begin{graphics}{0.8}{0.5}{0}{100}{-10}{10}{ve}{vs}{Titre}
\addPointsFromCSV{raw.csv}
\end{graphics}
@end_example
@end
\end{comment}

\newenvironment{graphicFigure}[9]
{
    \begin{center}
        \makeatletter
        \def\@captype{figure}
        \makeatother

        \newcommand{\TitleGraphic}{#9}%use var to print title 
        \begin{tikzpicture}
        \setGraphic %command to display with frenchb babel
    %\shorthandoff{:;!?} %Compulsory if frenchb package is used (from babel)
    \begin{axis}[width=#1\linewidth,height=#2\linewidth,xmin=#3,xmax=#4,  ymin=#5, ymax=#6, scale only axis,xlabel=#7,ylabel=#8] %grid=both
    }
    { 
    \end{axis}
        \end{tikzpicture}
        \captionof{figure}{\TitleGraphic}
    \end{center}
}


\newenvironment{graphic}[9]
{
        %\newcommand{\TitleGraphic}{#9}%use var to print title 
        \begin{tikzpicture}
        \setGraphic %command to display with frenchb babel
    \begin{axis}[width=#1\linewidth,height=#2\linewidth,xmin=#3,xmax=#4,  ymin=#5, ymax=#6, scale only axis,xlabel=#7,ylabel=#8, title=#9] %grid=both
    }
    { 
    \end{axis}
        \end{tikzpicture}
}
    
%############################################################
\begin{comment}
@begin
@command plot
@des 
Cette commande permet de tracer des surfaces 3D
@sed
@input Iitre du graphique
@input Equation à deux variables (x et y)
@begin_example 
\plot{Titre}{x*x+y*5}
@end_example
@end
\end{comment}   
    
\newcommand{\plot}[2]{
\begin{tikzpicture}
\setGraphic %command to display with frenchb babel
 \begin{axis}[title={#1}, xlabel=x, ylabel=y]
 \addplot3[surf,domain=0:360,samples=50]
 {#2};
 \end{axis}
 \end{tikzpicture}
 }
    
%############################################################
\begin{comment}
@begin
@command addPoints
@des 
Cette commande permet de tracer une courbe en passant une liste de coordonnées de points : (x1,y1)(x2,y2)(...)
Il faut que cette commande soit utilisées dans l'environnement \textbf{graphics}.
@sed
@input Couleur de la courbe
@input Liste des coordonnées des points. Chaque point est entre parenthèse et les coordonnées sont séparées par une virgule
@begin_example 
\addPoints{red}{(0,0)(5,0)(5,5)(10,5)(10,0)}
@end_example
@end
\end{comment}   

\newcommand{\addPoints}[2]{
    \addplot+[thick,mark=none, color=#1] coordinates{#2};
}

%############################################################
\begin{comment}
@begin
@command addTrace
@des 
Cette commande permet de tracer une courbe en passant une équation en fonction de x
Il faut que cette commande soit utilisées dans l'environnement \textbf{graphics}.
@sed
@input Couleur de la courbe
@input Début du domaine de définition
@input Fin du domaine de définition
@input Équation à tracer
@begin_example 
\addTrace{blue}{4}{8}{8x^2}
@end_example
@end
\end{comment}

\newcommand{\addTrace}[4]{
    \addplot [#1, domain=#2:#3, samples=400] {#4};
}

%############################################################
\begin{comment}
@begin
@command addPointsFromCSV
@des 
Cette commande permet de tracer une courbe en passant un fichier de données CSV (ou format TXT)
Il faut que cette commande soit utilisées dans l'environnement \textbf{graphics}.
@sed
@input Couleur de la courbe
@input délimitateur des données (comma=virgule, semicolon=point-vigule)
@input Fichier des données
@begin_example 
\addPointsFromCSV{blue}{semicolon}{raw.csv}
@end_example
@end
\end{comment}

\newcommand{\addPointsFromCSV}[3]{
    \IfFileExists{#3}{
    \addplot+[thick, mark=none, color=#1] table[mark=none,col sep=#2] {#3};
    }
    {\raiseError{[import failed]'#3' \stop}
    }
}

%############################################################
\begin{comment}
@begin
@command addLegend
@des 
Cette commande permet d'ajouter une légende au graphique
Il faut que cette commande soit utilisées dans l'environnement \textbf{graphics}.
@sed
@input Légendes de chaque courbe séparées par une virgule
@begin_example 
\addLegend{sig1, sig2, sig3}
@end_example
@end
\end{comment}

\newcommand{\addLegend}[1]{
    \legend{#1}
}

%############################################################%############################################################
%###### Package 'Layout' 
%###### This package contains some tools to set page layout or text
%###### Author  : Nicolas LE GUERROUE
%###### Contact : nicolasleguerroue@gmail.com
%############################################################
\RequirePackage{lmodern}
\RequirePackage[T1]{fontenc}
\RequirePackage[utf8]{inputenc}     %UTF-8 encoding
\RequirePackage{graphicx}           %Images
\RequirePackage{caption}            %légende
\RequirePackage{textcomp}           %special characters
\RequirePackage{fancyhdr}           %headers & footers
\RequirePackage{lastpage}           %page counter

\RequirePackage{float}              %image floating
\RequirePackage{wrapfig}
\RequirePackage{subcaption}         %Subcaption

\RequirePackage{geometry}

%############################################################

\begin{comment}
@begin
@command \setHeader
@des 
Cette commande permet de créer une page de garde minimaliste
@sed
@input Titre du document
@input Auteur(s) - Les retours à la ligne se font en utilisant la commande \\
@input La date
@begin_example 
\setHeader{Titre}{Auteur 1 \\ Auteur 2}{XX/XX/XXXX}
@end_example
@end
\end{comment}

\newcommand{\setHeader}[3]{
\title{#1}
\author{#2}
\date{#3}
\maketitle
}

% \makeatletter%
% \@ifclassloaded{babel}%
%   {BABEL is LOADED}%
%   {NOT}%
% \makeatother%


%###############################################################
\begin{comment}
@begin
@command \partImg
@des 
Cette commande permet de créer une page de partie avec une image
@sed
@input Titre de la partie
@input Source de l'image \\
@input Ratio
@begin_example 
\partImg{Partie}{Images/file.png}{0.2}
@end_example
@end
\end{comment}
\newcommand{\partImg}[3]{
    \part[#1]{#1 \\\vspace*{2cm} \makebox{\centering \includegraphics[width=#3\textwidth]{#2}}}
}


%###############################################################
\begin{comment}
@begin
@command \setHeaderImage
@des 
Cette commande permet de créer une page de garde avec une image centrée, un titre, sous titre en plus
@sed
@input Titre du document
@input Auteur(s) - Les retours à la ligne se font en utilisant la commande \\
@input La date
@begin_example 
\setHeaderImage{Titre}{Auteur 1 \\ Auteur 2}{XX/XX/XXXX}
@end_example
@end
\end{comment}

\newcommand{\setHeaderImage}[6]{
\begin{titlepage}
  \begin{sffamily}
  \begin{center}
    \includegraphics[scale=#2]{#1} \sn \sn
    \hfill
%\HRule \\[0.4cm]
\begin{center}
    {\Huge \textbf{#3}} \sn
    \textbf{#4}\sn \sn
\end{center}
\sn \sn
 #5 \sn
   \vfill
   {\large #6}
  \end{center}
  \end{sffamily}
\end{titlepage}
}


%### Définition du style de page 'classic' si report

\newcommand{\addPresentation}[6]{
\fancypagestyle{classic}{
    \rhead{#3}  
    \lhead{#1}
    \chead{#2}
    \rfoot{#6}  %Page courante / Nombre de page
    \cfoot{#5}
    \lfoot{#4}
}

\@ifclassloaded{report}{
\makeatletter
\renewcommand\chapter{\if@openright\cleardoublepage\else\clearpage\fi
                      \thispagestyle{classic} %Thème 'classic'
                      \global\@topnum\z@
                      \@afterindentfalse
                      \secdef\@chapter\@schapter}
\makeatother
}%End renew chapter

\@ifclassloaded{book}{
\makeatletter
\renewcommand\chapter{\if@openright\cleardoublepage\else\clearpage\fi
                      \thispagestyle{classic} %Thème 'classic'
                      \global\@topnum\z@
                      \@afterindentfalse
                      \secdef\@chapter\@schapter}
\makeatother
}%End renew chapter

\pagestyle{classic}
}


\newcommand{\setRightHeader}[1]{\rhead{#1}}
\newcommand{\setCenterHeader}[1]{\chead{#1}}
\newcommand{\setLeftHeader}[1]{\lhead{#1}}

\newcommand{\setRightFooter}[1]{\rfoot{#1}}
\newcommand{\setCenterFooter}[1]{\cfoot{#1}}
\newcommand{\setLeftFooter}[1]{\lfoot{#1}}



\newcommand{\setHeaderLine}[1]{ 
\renewcommand{\headrulewidth}{#1pt} 
}
\newcommand{\setFooterLine}[1]{ 
\renewcommand{\footrulewidth}{#1pt} 
}
  
\newcommand{\currentChapter}{\leftmark}

%Raccourcis
\@ifclassloaded{report)}{

\newcommand{\setAliasChapter}[1]{
\makeatletter
\renewcommand{\@chapapp}{#1}   %Le mot 'Chapitre' est remplacé par 'Section'
\makeatother
}
}%End if
{
  \newcommand{\setAliasChapter}[1]{
}
}

\@ifclassloaded{book)}{
\newcommand{\currentChapter}{\leftmark}
\newcommand{\setAliasChapter}[1]{
\makeatletter
\renewcommand{\@chapapp}{#1}   %Le mot 'Chapitre' est remplacé par 'Section'
\makeatother
}
}%End if


\newcommand{\currentPage}{\thepage/\pageref{LastPage}}%############################################################
%###### Package 'name' 
%###### This package contains ...
%###### Author  : Nicolas LE GUERROUE
%###### Contact : nicolasleguerroue@gmail.com
%############################################################
\RequirePackage{lmodern}
\RequirePackage[T1]{fontenc}
\RequirePackage[utf8]{inputenc}     %UTF-8 encoding
\RequirePackage{graphicx}           %Images
\RequirePackage{caption}            %légende
\RequirePackage{float}              %image floating
\RequirePackage{wrapfig}
\RequirePackage{subcaption}         %Subcaption
%################################################################
\begin{comment}
    @begin
    @command \img
    @des 
    Cette commande permet de mettre une image centrée avec un titre et une taille
    @sed
    @input Source et nom de l'image
    @input Titre de l'image
    @input Taille de l'image ]0;1+]
    @begin_example 
    \img{Image.png}{Titre de l'image}{0.5}
    @end_example
    @end
  \end{comment}
    
    \newcommand{\img}[3]{\IfFileExists{#1}{\begin{figure}[H]\centering\ \includegraphics[scale=#3]{#1}\caption{#2}\end{figure}}{\raiseWarning{[import failed]'#1'}}}

    %\newcommand{\imgd}[3]{\IfFileExists{#1}{\begin{figure}[H]\centering\ \includegraphics[scale=#3]{#1}\caption{#2}\end{figure}}{\raiseWarning{[import failed]'#1'}}}
    
    %################################################################
    \begin{comment}
    @begin
    @command \imgf
    @des 
    Cette commande permet de mettre une image centrée avec un titre et une taille de manière flottante avec le contenu
    @sed
    @input Source et nom de l'image
    @input Titre de l'image
    @input Taille de l'image ]0;1+] par rapport au texte
    @input Taille de l'image par rapport à l'espace qui lui est reservé
    
    @begin_example 
    \img{Image.png}{Titre de l'image}{0.5}{0.5}
    @end_example
    @end
    \end{comment}
    
    \newcommand{\imgf}[4]{
    \begin{subfigure}{#4\textwidth}
        \includegraphics[scale=#3]{#1}
      %\caption{#2}
    \end{subfigure}
    }
    %################################################################%############################################################
%###### Package 'name' 
%###### This package contains ...
%###### Author  : Nicolas LE GUERROUE
%###### Contact : nicolasleguerroue@gmail.com
%############################################################
\RequirePackage{enumitem}
\RequirePackage{pifont}

%################################################################

\newcommand{\setTriangleList}{\renewcommand{\labelitemi}{$\blacktriangleright$}}
\newcommand{\setCircleList}{\renewcommand{\labelitemi}{$\circ$}}
\newcommand{\setBulletList}{\renewcommand{\labelitemi}{$\bullet$}}
\newcommand{\setDiamondList}{\renewcommand{\labelitemi}{$\diamond$}}

\newcommand{\Triangle}{$\blacktriangleright$}
\newcommand{\Circle}{$\circ$}
\newcommand{\Bullet}{$\bullet$}
%\newcommand{\Diamond}{$\diamond$}


\newenvironment{items}[2]
{      
        \begin{itemize}[font=\color{#1}, label=#2]
           
    }
    { 
        \end{itemize}

}%############################################################
%###### Package 'name' 
%###### This package contains ...
%###### Author  : Nicolas LE GUERROUE
%###### Contact : nicolasleguerroue@gmail.com
%############################################################
\tcbuselibrary{listings,breakable, skins}
%\RequirePackage{amsmath}

%################################################################

\newtcbox{\file}{enhanced, nobeforeafter,tcbox raise base,boxrule=0.4pt,top=0mm,bottom=0mm,
  right=0mm,left=4mm,arc=1pt,boxsep=2pt,before upper={\vphantom{dlg}},
  colframe=orange!50!black,coltext=orange!25!black,colback=orange!10!white,
  overlay={\begin{tcbclipinterior}\fill[orange!75!white] (frame.south west)
    rectangle node[text=white,font=\sffamily\bfseries\tiny,rotate=90] {FILE} ([xshift=4mm]frame.north west);\end{tcbclipinterior}}}

\newtcbox{\lib}{enhanced, nobeforeafter,tcbox raise base,boxrule=0.4pt,top=0mm,bottom=0mm,
  right=0mm,left=4mm,arc=1pt,boxsep=2pt,before upper={\vphantom{dlg}},
  colframe=green!50!black,coltext=green!25!black,colback=green!10!white,
  overlay={\begin{tcbclipinterior}\fill[green!75!green!50!white] (frame.south west)
    rectangle node[text=white,font=\sffamily\bfseries\tiny,rotate=90] {LIB} ([xshift=4mm]frame.north west);\end{tcbclipinterior}}}

\newtcbox{\dir}{enhanced, nobeforeafter,tcbox raise base,boxrule=0.4pt,top=0mm,bottom=0mm,
  right=0mm,left=4mm,arc=1pt,boxsep=2pt,before upper={\vphantom{dlg}},
  colframe=blue!50!black,coltext=blue!25!black,colback=blue!10!white,
  overlay={\begin{tcbclipinterior}\fill[blue!75!blue!50!white] (frame.south west)
    rectangle node[text=white,font=\sffamily\bfseries\tiny,rotate=90] {DIR} ([xshift=4mm]frame.north west);\end{tcbclipinterior}}}


\newtcbox{\shortcut}{enhanced, nobeforeafter,tcbox raise base,boxrule=0.4pt,top=0mm,bottom=0mm,
  right=0mm,left=4mm,arc=1pt,boxsep=2pt,before upper={\vphantom{dlg}},
  colframe=purple!50!black,coltext=purple!25!black,colback=purple!10!white,
  overlay={\begin{tcbclipinterior}\fill[purple!75!blue!50!white] (frame.south west)
    rectangle node[text=white,font=\sffamily\bfseries\tiny,rotate=90] {KEY} ([xshift=4mm]frame.north west);\end{tcbclipinterior}}}

\newtcbox{\link}{enhanced, nobeforeafter,tcbox raise base,boxrule=0.4pt,top=0mm,bottom=0mm,
  right=0mm,left=4mm,arc=1pt,boxsep=2pt,before upper={\vphantom{dlg}},
  colframe=purple!50!black,coltext=purple!25!black,colback=purple!10!white,
  overlay={\begin{tcbclipinterior}\fill[purple!75!blue!50!white] (frame.south west)
    rectangle node[text=white,font=\sffamily\bfseries\tiny,rotate=90] {URL} ([xshift=4mm]frame.north west);\end{tcbclipinterior}}}


\newtcbox{\bin}{enhanced, nobeforeafter,tcbox raise base,boxrule=0.4pt,top=0mm,bottom=0mm,
  right=0mm,left=4mm,arc=1pt,boxsep=2pt,before upper={\vphantom{dlg}},
  colframe=blue!50!black,coltext=blue!25!black,colback=blue!10!white,
  overlay={\begin{tcbclipinterior}\fill[blue!75!blue!50!white] (frame.south west)
    rectangle node[text=white,font=\sffamily\bfseries\tiny,rotate=90] {BIN} ([xshift=4mm]frame.north west);\end{tcbclipinterior}}}

\newtcbox{\hexa}{enhanced, nobeforeafter,tcbox raise base,boxrule=0.4pt,top=0mm,bottom=0mm,
  right=0mm,left=4mm,arc=1pt,boxsep=2pt,before upper={\vphantom{dlg}},
  colframe=blue!50!black,coltext=blue!25!black,colback=blue!10!white,
  overlay={\begin{tcbclipinterior}\fill[blue!75!blue!50!white] (frame.south west)
    rectangle node[text=white,font=\sffamily\bfseries\tiny,rotate=90] {HEX} ([xshift=4mm]frame.north west);\end{tcbclipinterior}}}



\newtcbox{\ugrey}{nobeforeafter,colframe=gray,colback=gray!10!white,boxrule=0.5pt,arc=2pt,
  boxsep=0pt,left=2pt,right=2pt,top=2pt,bottom=2pt,tcbox raise base}


\newtcbox{\ugreen}{nobeforeafter,colframe=white,colback=green!40!white,boxrule=0pt,arc=0pt,
  boxsep=0pt,left=1pt,right=1pt,top=1pt,bottom=1pt,tcbox raise base}

\newtcbox{\ublue}{nobeforeafter,colframe=white,colback=blue!10!white,boxrule=0pt,arc=0pt,
  boxsep=0pt,left=1pt,right=1pt,top=1pt,bottom=1pt,tcbox raise base}

\newtcbox{\uorange}{nobeforeafter,colframe=white,colback=orange!40!white,boxrule=0pt,arc=0pt,
  boxsep=0pt,left=1pt,right=1pt,top=1pt,bottom=1pt,tcbox raise base}

\newtcbox{\ured}{nobeforeafter,colframe=white,colback=orange!40!red,boxrule=0pt,arc=0pt,
  boxsep=0pt,left=1pt,right=1pt,top=1pt,bottom=1pt,tcbox raise base}

\newtcbox{\highlight}[1]{nobeforeafter,colframe=white,colback=#1,boxrule=0pt,arc=0pt,
  boxsep=0pt,left=1pt,right=1pt,top=1pt,bottom=1pt,tcbox raise base}%############################################################
%###### Package 'Layout' 
%###### This package contains some tools to set page layout or text
%###### Author  : Nicolas LE GUERROUE
%###### Contact : nicolasleguerroue@gmail.com
%############################################################
\RequirePackage{lmodern}
\RequirePackage[T1]{fontenc}
\RequirePackage[utf8]{inputenc}     %UTF-8 encoding
\RequirePackage{graphicx}           %Images
\RequirePackage{caption}            %légende
\RequirePackage{textcomp}           %special characters

\RequirePackage{lastpage}           %page counter
%\RequirePackage{fullpage}          %entire page
\RequirePackage{float}              %image floating
\RequirePackage{wrapfig}
\RequirePackage{subcaption}         %Subcaption
\RequirePackage{listings}           %new env such as Bash, Python
%\RequirePackage{tcolorbox}          %messageBox
\RequirePackage[many]{tcolorbox}
\RequirePackage{color}              %colors
\RequirePackage{geometry}
\RequirePackage{listings}           %new env such as Bash, Python
\RequirePackage{xcolor}             %define new colors
\RequirePackage{xparse}
\RequirePackage{amssymb,amsthm}     %math
\RequirePackage{multirow}
\RequirePackage{tabularx}           %use for csv file [table]

\RequirePackage{lipsum}  
\RequirePackage{fontawesome}

%################################################################
\begin{comment}
@begin
@command \bold
@des 
Cette commande permet de mettre le texte en gras
@sed
@input Texte à mettre en gras
@begin_example 
\bold{texte en gras}
@end_example
@end
\end{comment}

\renewcommand{\bold}[1]{\textbf{#1}}

%################################################################
\begin{comment}
@begin
@command \italic
@des 
Cette commande permet de mettre le texte en italique
@sed
@input Texte à mettre en italique
@begin_example 
\italic{texte en italique}
@end_example
@end
\end{comment}

\newcommand{\italic}[1]{\textit{#1}}

%################################################################
\begin{comment}
@begin
@command \ib
@des 
Cette commande permet de mettre le texte en italique et en gras
@sed
@input Texte à mettre en italique et gras
@begin_example 
\ib{texte en italique et gras}
@end_example
@end
\end{comment}

\newcommand{\ib}[1]{\textit{\textbf{#1}}}

%################################################################
\begin{comment}
@begin
@command \bi
@des 
Cette commande permet de mettre le texte en italique et en gras
@sed
@input Texte à mettre en italique et gras
@begin_example 
\bi{texte en italique et gras}
@end_example
@end
\end{comment}

\newcommand{\bi}[1]{\textit{\textbf{#1}}}


\begin{comment}
@begin
@command \n
@des 
Cette commande permet de faire un saut de ligne
@sed
@begin_example 
Text \n Next text
@end_example
@end
\end{comment}

\newcommand{\n}{\\}

%################################################################
\begin{comment}
@begin
@command \sn
@des 
Cette commande permet de faire un espace vertical
@sed
@begin_example 
Text \sn Next text
@end_example
@end
\end{comment}

\newcommand{\sn}{\vskip 0.5cm}

%################################################################


\newtheorem{question}{Question}
\newtheorem{reponse}{$>>>$}
\newtheorem{propriete}{Propriété}
\newtheorem{proposition}{Proposition}
\newtheorem{remarque}{Remarque}
\newtheorem{exemple}{Exemple}
\newtheorem{definition}{Definition}


%#######################################
\newcommand{\messageBox}[5]{
% 1: title
% 2: color frame
% 3: color bg
% 4: content
% 5: titlecolor
%\messageBox{Remarque}{green}{white}{Erreur 0xff58}{black}
%\messageBox{Installation}{gray}{white}{DE}{black}
%\messageBox{Avertissement}{orange}{white}{dede}{black}
%\messageBox{Erreur}{red}{white}{dede}{black}
 \begin{tcolorbox}[title=#1,
colframe=#2!80,
colback=#3!10,
coltitle=#5!100,  
]
#4
\end{tcolorbox}
}



%######################### PAGE STYLE



\newcommand{\setFullPage}{\RequirePackage{fullpage}}%############################################################
%###### Package 'name' 
%###### This package contains tools to use links, url and other
%###### Author  : Nicolas LE GUERROUE
%###### Contact : nicolasleguerroue@gmail.com
%############################################################
\RequirePackage{hyperref}           %Url


%################################################################
\begin{comment}
@begin
@command \setParameters
@des 
Cette commande permet de définir les propriétés du document PDF (auteur, titre...)
@sed
@input Titre du PDF
@input Auteur(s)
@input Sujet du fichier PDF (courte phrase)
@input Créateur du fichier PDF
@input Mots-clés (liste)
@input Couleurs des liens
@input Couleurs des citations dans la bibliographie
@input Couleurs des liens de fichier
@input Couleurs des liens externe
@begin_example 
\setParameters {Tutoriel Latex} {Nicolas LE GUERROUE} {Tutoriel Latex pour la mise en place des outils} {\author}{Latex}{green}{blue}{blue}
@end_example
@end
\end{comment}

\newcommand{\setParameters}[8]{
\hypersetup{
    bookmarks=true,         % show bookmarks bar?
    unicode=true,          % non-Latin characters in Acrobat’s bookmarks
    pdftoolbar=true,        % show Acrobat’s toolbar?
    pdfmenubar=true,        % show Acrobat’s menu?
    pdffitwindow=false,     % window fit to page when opened
    pdfstartview={1},    % fits the width of the page to the window
    pdftitle={#1},    % title
    pdfauthor={#2},     % author
    pdfsubject={#3},   % subject of the document
    pdfcreator={#4},   % creator of the document
    pdfproducer={#4}, % producer of the document
    pdfkeywords={#5}, % list of keywords
    pdfnewwindow=true,      % links in new PDF window
    colorlinks=true,       % false: boxed links; true: colored links
    linkcolor=black,          % color of internal links (change box color with linkbordercolor)
    citecolor=#6,        % color of links to bibliography
    filecolor=#7,         % color of file links
    urlcolor=#8        % color of external links
}}%############################################################
%###### Package 'Maths' 
%###### This package contains some tools to set mathematic tools
%###### Author  : Nicolas LE GUERROUE
%###### Contact : nicolasleguerroue@gmail.com
%############################################################
%Vecteur à 3 composantes  
\newcommand{\evec}[3]{\left (\begin{array}{ccc} #1 \\ #2 \\ #3\end{array} \right )}
%Matrice à 3 composantes
\newcommand{\emat}[3]{\left [\begin{array}{ccc} #1 \\ #2 \\ #3\end{array} \right ]}
% \emat{a & b & c}{d & e & f}{g & h & i}%############################################################
%###### Package 'Titles' 
%###### This package contains some tools to set title color
%###### Author  : Nicolas LE GUERROUE
%###### Contact : nicolasleguerroue@gmail.com
%############################################################
\RequirePackage{pdfpages} 
%################################################################

%\includepdf[⟨options⟩]{⟨fichier⟩}\RequirePackage{color}              %colors
\RequirePackage{listings}           %new env such as Bash, Python
\RequirePackage{xcolor}             %define new colors
\RequirePackage{graphicx}             %define new colors
\RequirePackage{xparse}
\RequirePackage{amssymb,amsthm}     %math
\RequirePackage{multirow}
\RequirePackage{tabularx}           %use for csv file [table]
\RequirePackage{textcomp}           %special characters

%###############################################################
%### Couleurs programmation
\definecolor{green}{rgb}{0,0.6,0}
\definecolor{gray}{rgb}{0.5,0.5,0.5}
\definecolor{purple}{rgb}{0.58,0,0.82}
\definecolor{bg}{rgb}{0.98,0.98,0.98}
%###############################################################
%### définition du style 'Python' [begin{Python}]
\lstdefinestyle{Python}{
    backgroundcolor=\color{bg},   
    commentstyle=\color{green},
    keywordstyle=\color{blue},
    numberstyle=\tiny\color{gray},
    stringstyle=\color{purple},
    basicstyle=\ttfamily\footnotesize,
    breakatwhitespace=false,         
    breaklines=true,                 
    captionpos=b,                    
    keepspaces=true,                 
    numbers=none,   %left, right, none 
    numbersep=5pt,                  
    showspaces=false,                
    showstringspaces=false,
    showtabs=false,  
    tabsize=2
}
\lstset{style=Python}
\lstset{literate=
  {á}{{\'a}}1 {é}{{\'e}}1 {í}{{\'i}}1 {ó}{{\'o}}1 {ú}{{\'u}}1
  {Á}{{\'A}}1 {É}{{\'E}}1 {Í}{{\'I}}1 {Ó}{{\'O}}1 {Ú}{{\'U}}1
  {à}{{\`a}}1 {è}{{\`e}}1 {ì}{{\`i}}1 {ò}{{\`o}}1 {ù}{{\`u}}1
  {À}{{\`A}}1 {È}{{\'E}}1 {Ì}{{\`I}}1 {Ò}{{\`O}}1 {Ù}{{\`U}}1
  {ä}{{\"a}}1 {ë}{{\"e}}1 {ï}{{\"i}}1 {ö}{{\"o}}1 {ü}{{\"u}}1
  {Ä}{{\"A}}1 {Ë}{{\"E}}1 {Ï}{{\"I}}1 {Ö}{{\"O}}1 {Ü}{{\"U}}1
  {â}{{\^a}}1 {ê}{{\^e}}1 {î}{{\^i}}1 {ô}{{\^o}}1 {û}{{\^u}}1
  {Â}{{\^A}}1 {Ê}{{\^E}}1 {Î}{{\^I}}1 {Ô}{{\^O}}1 {Û}{{\^U}}1
  {Ã}{{\~A}}1 {ã}{{\~a}}1 {Õ}{{\~O}}1 {õ}{{\~o}}1 {’}{{'}}1
  {œ}{{\oe}}1 {Œ}{{\OE}}1 {æ}{{\ae}}1 {Æ}{{\AE}}1 {ß}{{\ss}}1
  {ű}{{\H{u}}}1 {Ű}{{\H{U}}}1 {ő}{{\H{o}}}1 {Ő}{{\H{O}}}1
  {ç}{{\c c}}1 {Ç}{{\c C}}1 {ø}{{\o}}1 {å}{{\r a}}1 {Å}{{\r A}}1
  {€}{{\euro}}1 {£}{{\pounds}}1 {«}{{\guillemotleft}}1
  {»}{{\guillemotright}}1 {ñ}{{\~n}}1 {Ñ}{{\~N}}1 {¿}{{?`}}1
}
\lstnewenvironment{Python}[1]{\lstset{language=Python, title=#1}}{}
%###############################################################
%### Couleurs programmation
\definecolor{green}{rgb}{0,0.6,0}
\definecolor{gray}{rgb}{0.5,0.5,0.5}
\definecolor{purple}{rgb}{0.58,0,0.82}
\definecolor{bg}{rgb}{0.98,0.98,0.98}
%###############################################################
%### définition du style 'Cpp' [begin{Cpp}]
\lstdefinestyle{Cpp}{
    basicstyle=\ttfamily,
    columns=fullflexible,
    keepspaces=true,
    upquote=true,
    showstringspaces=false,
    commentstyle=\color{olive},
    keywordstyle=\color{blue},
    identifierstyle=\color{violet},
    stringstyle=\color{purple},
    numbers=none,
    language=c,
}
\lstnewenvironment{Cpp}[1]{\lstset{style=Cpp, title=#1}}{}

%###############################################################
%### définition du style 'Bash' [begin{Bash}]
\lstdefinestyle{Bash}{
    basicstyle=\ttfamily,
    columns=fullflexible,
    keepspaces=true,
    upquote=true,
    frame=trBL, %frame=trBL, shadowbox, trrb
    %frameround=fttt,
    numbers=none,
    showstringspaces=false,
    commentstyle=\color{olive},
    keywordstyle=\color{blue},
    identifierstyle=\color{black},
    stringstyle=\color{gray},
    language=bash,
    morekeywords={sudo,apt-get,install, autoremove, update, upgrade}
}
\lstnewenvironment{Bash}[1]{\lstset{style=Bash, title=#1}}{}%############################################################
%###### Package 'Titles' 
%###### This package contains some tools to set title color
%###### Author  : Nicolas LE GUERROUE
%###### Contact : nicolasleguerroue@gmail.com
%############################################################
%\RequirePackage{titlesec} 
\RequirePackage[explicit]{titlesec} 
%################################################################

\@ifclassloaded{report}{
    \titleformat{\chapter}[display] {\fontsize{17pt}{12pt}\selectfont \bfseries}{\textcolor{red} {\chaptertitlename\ \thechapter: #1}}{20pt}{\Huge}
}

\titleformat{\section}[display] {\fontsize{17pt}{12pt}\selectfont \bfseries}{\textcolor{DarkBlue} {#1}}{20pt}{\Huge}
\titleformat{\subsection}[display] {\fontsize{15pt}{12pt}\selectfont \bfseries}{\textcolor{orange} {#1}}{20pt}{\large}
\titlespacing*{\section}{0pt}{20pt}{-30pt}
\titlespacing*{\subsection}{0pt}{20pt}{-20pt}
%\usepackage[Lenny]{fncychap}   %Sonny, Lenny, Glenn, Conny, Rejne, Bjarne
%#############################################################
%#### Settings
%#############################################################
%#############################################################
%If you want to set on fullpage, change this bloc
%#############################################################
\geometry{hmargin=2cm,vmargin=2.5cm}
%#############################################################
%#############################################################
%If you want rename the chapter name, change the value of argument
%#############################################################
\setAliasChapter{Section}
%#############################################################
%#############################################################
%If you want to add presentation, modify the next bloc
%The firt line is about the header : 
% {left content}{center content}{right content}
%The second line is about the footer : 
% {left content}{center content}{right content}
%####
%to get the current chapter name, use \currentChapter command as content
%to get the current number page, use \currentPagecommand as content
\addPresentation
{Synthèse Latex} {} {\currentChapter}
{Bibliothèque Utils} {} {\currentPage}
%#############################################################
%Change the width of footer line and header line
%To delete it, set value to 0
\setHeaderLine{0.2}
\setFooterLine{0.2}

\setcounter{tocdepth}{2} %depth of table of content
\setcounter{secnumdepth}{2}

%#############################################################
%Change the name of the section "Nomenclature"
%To delete it, set value to 0
%\renewcommand{\nomname}{Conventions}
%Setting up the parameter of PDF file as name, author...
\begin{comment}
@input Titre du PDF
@input Auteur(s)
@input Sujet du fichier PDF (courte phrase)
@input Créateur du fichier PDF
@input Producteur du fichier PDF
@input Mots-clés (liste)
@input Couleurs des liens
@input Couleurs des citations dans la bibliographie
@input Couleurs des liens de fichier
\end{comment}
\setParameters {Tutoriel Latex} {Nicolas Le Guerroué} {Bibliothèque Utils} {Nicolas Le Guerroué}{Latex}{green}{blue}{blue}
  


\titleformat{\chapter}
  {\gdef\chapterlabel{}
   \normalfont\sffamily\Huge\bfseries\scshape}
  {\gdef\chapterlabel{\thechapter)\ }}{0pt}
  {\begin{tikzpicture}[remember picture,overlay]
    \node[yshift=-2cm] at (current page.north west)
      {\begin{tikzpicture}[remember picture, overlay]
        \draw[fill=MediumBlue] (0,0) rectangle
          (\paperwidth,2cm);
        \node[anchor=east,xshift=.9\paperwidth,rectangle,
              rounded corners=20pt,inner sep=11pt,
              fill=white]
              {\color{black}#1};%\chapterlabel
          \node[anchor=west,yshift=1cm,xshift=2cm,inner sep=11pt,
              fill=MediumBlue]
              {\color{white}{\large Synthèse Latex}};%\chapterlabel
       \end{tikzpicture}
      };
   \end{tikzpicture}
  }
  \titlespacing*{\chapter}{0pt}{50pt}{-60pt}
\begin{document}
\setHeader{A}{A}{A}
\tableofcontents
\setcounter{page}{2}
  
\chapter{La bataille de Waterloo}

\section{Introduction}

À première vue, La bataille de Waterloo voit la confrontation finale entre deux puissances. Or, la troisième puissance partipante est vite oubliée : il s'agit de la Prusse, nation ayant émergée sous le règne du roi Frédéric.

Dans l'histoire collective, le vainquu a éclipsé le vainqueur. Combien de personnes conaissent le nom du vainqueur de Waterloo ? Sans doute un anglais, mais qui ?

La légende napoléonnienne a fait le reste, seul le dernier carré de la garde et le mot de Cambronne sont restés.



\begin{items}{orange}{\Triangle}
    \item A
    \item B
    \item C
\end{items}

\subsubsection{Introduction}

eu, accumsan eleifend, sagittis quis, diam. Duis eget orci sit amet orci dignissim rutrum.
Nam dui ligula, fringilla a, euismod sodales, sollicitudin vel, wisi. Morbi auctor lorem
non justo. Nam lacus libero, pretium at, lobortis vitae, ultricies et, tellus. Donec aliquet,
tortor sed accumsan bibendum, erat ligula aliquet magna, vitae ornare odio metus a mi.
Morbi ac orci et nisl hendrerit mollis. Suspendisse ut massa. Cras nec ante. Pellentesque
a nulla. Cum sociis natoque penatibus et magnis dis parturient montes, nascetur ridicu-
lus mus. Aliquam tincidunt urna. Nulla ullamcorper vestibulum turpis. Pellentesque cur-
sus



\begin{figure}[!t]
  \centering
  \begin{subfigure}{.22\linewidth}
    \centering
   \begin{tikzpicture}
  \begin{axis}[height=1\linewidth]
    \addplot+[thick,mark=none, color=blue] coordinates{ (0,1)(1,1)(1.01,0)(3,0) };
  \end{axis}
  \end{tikzpicture}
    \caption{Passe-bas}
    \label{fig2:sub1}
  \end{subfigure}%
  \hfill
  \begin{subfigure}{.22\linewidth}
    \centering
     \begin{tikzpicture}
  \begin{axis}[height=1\linewidth]
    \addplot+[thick,mark=none, color=blue] coordinates{(0,0) (0.99,0) (1,1) (3,1)} ;
  \end{axis}
  \end{tikzpicture}
    \caption{Passe-haut}
    \label{fig2:sub2}
  \end{subfigure}\hfill
  \begin{subfigure}{.22\linewidth}
    \centering
     \begin{tikzpicture}
  \begin{axis}[height=1\linewidth]
    \addplot+[thick,mark=none, color=blue] coordinates{(0,0) (0.99,0) (1,1) (2,1) (2.01,0) (3,0)};
  \end{axis}
  \end{tikzpicture}
    \caption{Passe-bande}
    \label{fig2:sub3}
  \end{subfigure}\hfill
  \begin{subfigure}{.22\linewidth}
    \centering
     \begin{tikzpicture}
  \begin{axis}[height=1\linewidth]
    \addplot+[thick,mark=none, color=blue] coordinates{(0,1) (0.99,1) (1,0) (2,0) (2.01,1)(3,1)};
  \end{axis}
  \end{tikzpicture}
    \caption{Coupe bande}
    \label{fig2:sub4}
  \end{subfigure}
  \caption{Filtres en "mur de briques"}
  \label{fig:murbrique}
  \end{figure}

%\chapter{La Bilbiothèque Layout}
hello

  \begin{figure}[h]  
  \centering 
    \begin{subfigure}[b]{0.4\linewidth}
      \begin{graphic}{0.8}{1}{0}{1.2}{-1}{5}{t(s)}{Tension V}{h}
        \addPointsFromCSV{red}{comma}{src_examples/sinus.txt}
        \addLegend{sin(t)}
        \end{graphic}%NO END  LINE HERE
      \caption{No interaction} 
    \end{subfigure}
  \begin{subfigure}[b]{0.4\linewidth}
    \begin{graphic}{0.8}{1}{0}{4}{-0.3}{0.3}{t(s)}{Tension V}{g}
      \addPointsFromCSV{blue}{comma}{src_examples/jack01.txt}
      \addLegend{g(t)}
      \end{graphic}%NO END  LINE HERE
  \caption{Interaction}
  \end{subfigure}
  \caption{Les tensions de service}
  \end{figure}  

% Ce document a pour but de présenter les fonctionnalités de la bibliothèque Utils. \n
% Chaque bibliothèque doit être indépendante afin de fonctionner correctement.


%\chapter{Bibliothèque Layout}
% \addQuote{En temps de guerre, la vérité est si précieuse qu'elle doit etre escortée par une garde de mensonge}{Sir Winston Spencer Churchill,\\ Duc de Malborought}
% \addQuote{Si Dieu nous fait la grâce de perdre encore une pareille bataille, Votre Majesté peut compter que ses ennemis sont détruits}{Maréchal de Villars, \\ 1709}


% \section{Les commandes}

% \paragraph{la bibliothèque 'Layout' permet de gérer la disposition et la mise en forme de certains élements du document. \\}

% \subsection{Mise en forme de la page de garde avec image}

% \italic{Se référer au code source de la page.\footnote{Commande setHeaderImage}}
% \index{setHeaderImage}

% \subsection{Mise en forme de la page de garde sans image}

% \bi{Se référer à la commande setHeader}
% \index{setHeader}

% \subsection{Ajout d'un trait entre l'en-tête et le corps de la page}

% \bi{Se référer à la commande setHeaderLine}
% \index{setHeaderLine}

% \subsection{Ajout d'un trait entre le corps de la page et le bas de page}

% \bi{Se référer à la commande setFooterLine}
% \index{setFooterLine}

% \subsection{Définition de la présentation globale des pages}

% \bi{Se référer à la commande addPresentation}
% \index{setFooterLine}


% \subsection{Redéfinir les titres des chapitres}

% \bi{Se référer à la commande setAliasChapter}
% \index{setAliasChapter}

% \subsection{Mettre le document en pleine page}

% \bi{Se référer à la commande setFullPage}
% \index{setFullpage}

% \subsection{Mise en gras}

% \bold{Texte en gras}
% \index{bold}

% \subsection{Mise en italique}

% \italic{texte en italic}
% \index{italic}

% \subsection{Mise en gras et italique}

% \ib{Texte en gras et italique}
% \index{ib}

% \subsection{Ajout d'une image non-flottante}

% \img{\rootImages/tux.png}{Légende de l'image}{0.5}
% \index{img}

% \subsection{Ajout d'une image flottante}

% Lorem ipsum dolor sit amet, consectetuer adipiscing elit. Ut purus elit, vestibulum
% ut, placerat ac, adipiscing vitae, felis. Curabitur dictum gravida mauris. Nam arcu li-
% bero, nonummy eget, consectetuer id, vulputate a, magna. Donec vehicula augue eu
% neque. Pellentesque habitant morbi tristique senectus et netus et malesuada fames ac
% turpis egestas. Mauris ut leo. Cras viverra metus rhoncus sem. Nulla et lectus vestibu-
% lum urna fringilla ultrices. Phasellus eu tellus sit amet tortor gravida placerat. Integer
% sapien est, iaculis in, pretium quis, viverra ac, nunc. Praesent eget sem vel leo ultrices
% bibendum. Aenean faucibus. Morbi dolor nulla, malesuada eu, pulvinar at, mollis ac,
% nulla. Curabitur auctor semper nulla. Donec varius orci eget risus. Duis nibh mi, congue

% \imgf{\rootImages/tux.png}{Tux}{0.6}{0.5}
% \index{imgf}

% eu, accumsan eleifend, sagittis quis, diam. Duis eget orci sit amet orci dignissim rutrum.
% Nam dui ligula, fringilla a, euismod sodales, sollicitudin vel, wisi. Morbi auctor lorem
% non justo. Nam lacus libero, pretium at, lobortis vitae, ultricies et, tellus. Donec aliquet,
% tortor sed accumsan bibendum, erat ligula aliquet magna, vitae ornare odio metus a mi.
% Morbi ac orci et nisl hendrerit mollis. Suspendisse ut massa. Cras nec ante. Pellentesque
% a nulla. Cum sociis natoque penatibus et magnis dis parturient montes, nascetur ridicu-
% lus mus. Aliquam tincidunt urna. Nulla ullamcorper vestibulum turpis. Pellentesque cur-
% sus
% luctus
% mauris.

% \subsection{Ajout d'un espace vertical}

% Lorem ipsum dolor sit amet, consectetuer adipiscing elit. Ut purus elit, vestibulum
% ut, placerat ac, adipiscing vitae, felis. Curabitur dictum gravida mauris. Nam arcu li- \sn
% \index{sn}
% bero, nonummy eget, consectetuer id, vulputate a, magna. Donec vehicula augue eu
% neque. Pellentesque habitant morbi tristique senectus et netus et malesuada fames ac
% turpis egestas. Mauris ut leo. Cras viverra metus rhoncus sem. Nulla et lectus vestibu-

% \subsection{Ajout d'une boite de dialogue}

% \messageBox{Message}{orange}{white}{Voici un message}{black}
% \index{messageBox}
% \messageBox{Message}{red}{white}{Et un autre message}{black}
% \messageBox{Message}{green}{white}{bref...}{white}

% \section{Les environnements de la bibliothèque}

% \subsection{Environnement 'question' }

% \begin{question}
% Quelle heure est-il ?
% \end{question}
% \index{question}

% \subsection{Environnement 'reponse' }

% \begin{reponse}
% il est 18 h.
% \end{reponse}
% \index{reponse}

% \subsection{Environnement 'propriete' }

% \begin{propriete}
% Un produit scalaire est commutatif.
% \end{propriete}
% \index{propriete}

% \subsection{Environnement 'proposition' }

% \begin{proposition}
% Les chats sont des mammifères.
% \end{proposition}
% \index{proposition}

% \subsection{Environnement 'remarque' }

% \begin{remarque}
% Les chats sont des mammifères.
% \end{remarque}
% \index{remarque}

% \subsection{Environnement 'exemple' }

% \begin{exemple}
% Ceci est un exemple d'exemple
% \end{exemple}
% \index{exemple}

% \subsection{Environnement 'definition' }

% \begin{definition}
% Une phrase est un ensemble de mots.
% \end{definition}
% \index{definition}


%###########################################################################
%LINKS
%###########################################################################

% \chapter{Bilbiothèque 'Links'}


% \section{Les commandes}


% \subsection{définition des liens dans le document}

% \bi{Se référer à la commande setParameters}
% \index{setParameters}


%###########################################################################
%MATHS
%###########################################################################

% \chapter{Bilbiothèque 'maths'}


% \section{Les commandes}


% \subsection{Création d'une matrice 3*3}


% $\emat{a & b & c}{d & e & f}{g & h & i}  $
% \index{emat}

% \subsection{Création d'un vecteur à trois dimensions}
% $\evec{a}{b}{c}  $
% \index{evec}




%###########################################################################
%PROGRAMMING
%###########################################################################

% \chapter{Bilbiothèque 'Programming'}


% \section{Les environnements}


% \subsection{Affichage d'un code C/C++ avec titre}


% \begin{Cpp}{Titre}
% #include <iostream>

% #define CONST 1

% int var = 1;
% float 

% int main() {
  
%   call();
%   return 0;

% }//End main

% \end{Cpp}
% \index{cpp}
% \subsection{Affichage d'un code C/C++ sans titre}


% \begin{Cpp}
% #include <iostream>

% #define CONST 1 #const var

% int var = 1;
% float g = 2.5;
% ...

% \end{Cpp}


% \subsection{Affichage d'un code Python avec titre}


% \begin{Python}{Titre du code}
% def call(input):

%   """docstring"""
%   a = input
%   for elem in a:
%     print(elem) #show
% \end{Python}
% \index{python}
% \subsection{Affichage d'un code Python sans titre}


% \begin{Python}
% def call(input):

%   """docstring"""
%   ...
% \end{Python}


% \subsection{Affichage d'un code Bash avec titre}


% \begin{Bash}{Titre du code}
% sudo apt-get -y update
% sudo apt-get -y upgrade
% echo -e "content"
% \end{Bash}
% \index{bash}
% \subsection{Affichage d'un code bash sans titre}

% \bold{Ce type d'affichage n'est pas encore supporté par la bibliothèque.}




%###########################################################################
%PROGRAMMING
%###########################################################################

% \chapter{Bilbiothèque 'Labels'}

% \section{Les labels de surlignage}

% \subsection{Affichage d'un label de surlignage vert}
<true>
% \ugreen{content}
% \index{ugreen}
% \subsection{Affichage d'un label de surlignage orange}

% \uorange{content}
% \index{uorange}
% \subsection{Affichage d'un label de surlignage rouge}

% \ured{content}
% \index{ured}
% \subsection{Affichage d'un label de surlignage bleu}

% \ublue{content}
% \index{ublue}
% \subsection{Affichage d'un label de surlignage gris encadré}

% \ugrey{content}
% \index{ugrey}

% \section{Le surlignage générique}

% \highlight{orange}{Texte à surligner}

% Ce texte est généré par la fonction \bold{highlight}

% \section{Les labels avancés}

% \subsection{Affichage d'un label de bibliothèque}

% \lib{Gparted}
% \index{lib}
% \subsection{Affichage d'un label de fichier}

% \file{Fichier.txt}
% \index{file}
% \subsection{Affichage d'un label de dossier ou d'emplacement}

% \dir{/var/www/html}
% \index{dir}
% \subsection{Affichage d'un label de raccourci-clavier}

% \shortcut{Ctrl+'A'}
% \index{shortcut}

% \section{Les labels numérique}

% \subsection{Affichage d'un label de nombre hexadécimal}

% \hexa{FF45DE}
% \index{hexa}
% \subsection{Affichage d'un label de nombre binaire}

% \bin{0100101}
% \index{bin}




% \highlight{blue}{content is here}
% \highlight{green}{content is here}
% \highlight{pink}{content is here}

%###########################################################################
%GRAPHIC
%###########################################################################

% \chapter{Bilbiothèque 'Graphic'}

% \section{Graphiques 2D}
% \subsection{Affichage d'un graphique 2D avec insertion des données depuis un fichier txt (csv)}


% \begin{graphics}{0.8}{0.6}{0}{2.1}{-1.1}{1.1}{t(ms)}{vs}{Oscilloscope}
% \addPointsFromCSV{red}{comma}{src_examples/input_1.txt}
% \addPointsFromCSV{blue}{comma}{src_examples/input_2.txt}
% \addLegend{voie A, voie B}
% \end{graphics}
% \index{graphic}
% \index{addPointsFromCSV}

% \subsection{Affichage d'un graphique 2D avec insertion des données depuis une liste de points}


% \begin{graphics}{0.8}{0.4}{0}{40}{-1}{6}{t(s)}{Tension (V)}{Signal numérique}
% \addPoints{blue}{(0,0)(10,0)(10,5)(15,5)(15,0)(20,0)(20,5)(25,5)(25,0)(30,0)(30,5)(35,5)(35,0)(100,0)}
% \addLegend{Tension (V)}
% \end{graphics}


% \subsection{Affichage de plusieurs graphiques avec insertion des données depuis un fichier externe}

% \begin{figure}
%   \centering
%   \subfloat[fig 1]{
%     \begin{graphic}{0.9}{0.4}{0.0}{5}{-2}{2}{t(ms)}{vs}{$sin(t)$}
%     \addPointsFromCSV{red}{comma}{src_examples/sinxx.txt}
%     \addPointsFromCSV{blue}{comma}{src_examples/gauss.txt}
%     \addLegend{Sinus}
%     \end{graphic}} 
%   \subfloat[fig 2]{\includegraphics[width = 1in]{\rootImages/tux}}\\
%   \subfloat[fig 3]{\includegraphics[width = 1in]{\rootImages/tux}}
%   \subfloat[fig 4]{\includegraphics[width = 1in]{\rootImages/tux}} 
%   \caption{Add your own figures before compiling}
%   \label{some example}
%   \end{figure}


% \begin{figure}

%     \begin{subfigure}{.30\textwidth}
%      \begin{graphic}{0.9}{0.4}{0.0}{5}{-2}{2}{t(ms)}{vs}{$sin(t)$}
%     \addPointsFromCSV{red}{comma}{src_examples/sinxx.txt}
%     \addPointsFromCSV{blue}{comma}{src_examples/gauss.txt}
%     \addLegend{Sinus}
%     \end{graphic}

%     \begin{graphic}{0.9}{0.4}{0.0}{5}{-2}{2}{t(ms)}{vs}{$sin(t)$}
%       \addPointsFromCSV{red}{comma}{src_examples/gauss.txt}
%       \addLegend{Sinus}
%       \end{graphic}

%     \end{subfigure}
    

%     \begin{subfigure}{.30\textwidth}
%      \begin{graphic}{0.9}{0.4}{0.0}{5}{-2}{2}{t(ms)}{vs}{$sin(t)$}
%     \addPointsFromCSV{red}{comma}{src_examples/gauss.txt}
%     \addLegend{Sinus}
%     \end{graphic}

%     \end{subfigure}

%   \caption{plots of....}
%   \label{fig:fig}
%   \end{figure}

\newpage

% \subsection{Affichage d'un graphique 2D avec insertion des données depuis une équation}



% \begin{graphicFigure}{0.4}{0.4}{0}{3}{-1}{5}{t(s)}{Tension V}{Signal analogique}
% \addPointsFromCSV{red}{comma}{src_examples/input_2.txt}
% \addLegend{AA}
% \end{graphicFigure}


  
\begin{figure}[!h]
  \centering
\begin{tabular}{|c|c|c|c|}
  \hline
  $U_A$ (V) & $U_B$ (V) & Sens du courant & $U_A-U_B$\\
  \hline
  10 & 5 & \colors{blue}{De A vers B} & 5\\
  \hline
  5 & 10 & \colors{blue}{de B vers A} & -5\\
  \hline
  5 & 5 & \colors{blue}{Aucun courant ne circule} & 0\\
  \hline
\end{tabular}
\caption{Réponse sur le sens du courant en fonction des tensions $U_A$ et $U_B$}
\end{figure}



\begin{numeric}{exemple numérique}
  D1 &  20{C}   \\
  D2 &  [green] 1H1L1L1L1H1L1L1H1L1H1L1H1L1H1L1H  \\
  D3 &  [brown] 20{C}   \\
  D4 &  [red] 1H1L1L1L1H1L1L1H1L1H1L1H1L1H1L1H  \\
  D5 &  [blue] 20{C}   \\
  D6 &  [black] 1H1L1L1L1H1L1L1H1L1H1L1H1L1H1L1H  \\
  D7 &  [black] 1H1L1L1L1H1L1L1H1L1H1L1H1L1H1L1H  \\
  D8 & 8D5U7U5D \\
  D9 & LLL 2{0.1H 0.1L} 0.6H HH \\
  D10 & ZZ G ZZ G XX G X \\
  D11 & [d] 4{5D{Text}} 0.2D \\
  D12 & [L][timing/slope=1.0] HL HL HL HL HL \\
\end{numeric}

Notre signal d'horloge (\texttiming{[blue]CCCCCC}) provient d'un oscillateur à quartz.
\texttiming[timing/draw grid]{LHLHLHLHLHLHLHL}


\begin{numeric}{Chronogramme du compteur 4 bits}
  INPUT &  CC [blue]16{CC} CCC   \\
  D0 &  HL 8{LHHL} LHL   \\
  D1 &  H  4{LLLLHHHH} LLLL \\
  D2 &  H 2{LLLLLLLLHHHHHHHH} LLLL   \\
  D3 &  H{LLLLLLLLLLLLLLLLHHHHHHHHHHHHHHHH} LLLL  \\
  END &  LL [green]14{LL} LHHLLL  \\
  VALUE & L 2D{0} 2D{1} 2D{2} 2D{3} 2D{4} 2D{5} 2D{6} 2D{7} 2D{8} 2D{9} 2D{10} 2D{11} 2D{12} 2D{13} 2D{14} 2D{15} 2D{0} 2D{1}  \\
\end{numeric}%

\begin{schema} {Convertisseur Analogique-Numérique}

  \addPower{6,5}{power1}{$+5V$}
  \addGround{4,0}{gnd1}{}

  \setDeviceBackgroundColor{white}
  \setRotate{0}
  \addLogicGate{0,0}{mynor}{nor}{}{A}{B}{G1}

  \setDeviceBackgroundColor{green}
  \addLogicGate{0,2}{mynand}{nand}{}{C}{D}{G2}
  \addLogicGate{2,1}{myor}{or}{}{}{}{G3}
  \resetColors
          
  \addTransistor{6,1}{npnA}{nmos}{B}{C}{E}
  \addTransistor{6,3}{pnpA}{pmos}{b}{e}{c}

  \resetColors
  \addTransistor{10,2}{npnR}{nmos}{b}{e}{c}

  \addWire{mynor.out}{myor.in 2}{\orthogonalWireA}
  \addWire{mynand.out}{myor.in 1}{\orthogonalWireA}

  %\addWire{myor.out}{npnA.B}{\orthogonalWireA}
  \addWire{mynand.out}{pnpA.B}{\orthogonalWireA}
  \addWire{pnpA.C}{npnA.C}{\orthogonalWireA}

  \addWire{pnpA.E}{power1}{\orthogonalWireA}

  \addWire{npnA.E}{gnd1}{\orthogonalWireA}

  \addNode{$(pnpA.C)+(1,0)$}{node1}{}
  \addWire{pnpA.C}{node1}{\orthogonalWireA}

  \setDeviceBackgroundColor{red}
  \addLed{myor.out}{\Right}{npnA.B}{\orthogonalWireA}{L1}
  \addResistor{node1}{\Right}{npnR.B}{\orthogonalWireA}
          
\end{schema}


\begin{schema}{Pont de Graëtz}

  \addNode{0,0}{a}{A}
  \addNode{2.5,2.5}{c}{C}
  \addNode{5,5}{b}{B}
  \addLogicGate{0,2}{mynand}{nand}{}{C}{D}

  \addTransistor{4,1}{npnA}{nmos}{B}{C}{E}

  \addResistor{mynand.out}{1.5,0}{b}{\orthogonalWireA}

  \addPower{1,4}{power1}{$+5V$}
  \addResistor{power1}{\Down}{npnA.B}{\orthogonalWireB}

\end{schema}


% \begin{graphics}{0.8}{0.4}{0}{40}{-1}{6}{t(s)}{Tension (V)}{Signal numérique}
% \addPoints{blue}{(0,0)(10,0)(10,5)(15,5)(15,0)(20,0)(20,5)(25,5)(25,0)(30,0)(30,5)(35,5)(35,0)(100,0)}
% \addLegend{Tension (V)}
% \end{graphics}

\begin{tabular}{|c|c||c|}
          \hline
          $A$& $B$&$S$\\
          \hline
          0 & 0 & 0\\
          \hline
          0 & 1 & 1\\
          \hline
          1 & 0 & 1\\
          \hline
          1 & 1 & 1\\
          \hline
\end{tabular}

\begin{figure}
\begin{subfigure}{.5\textwidth}
\centering
% include first image
\centering
\begin{schema}{Pont de Graëtz}
\setBipolesLength{0.5}
\addNode{0,0}{a}{A}
\addNode{2.5,2.5}{c}{C}
\addNode{5,5}{b}{B}
\addLogicGate{0,2}{mynand}{nand}{}{C}{D}

\addTransistor{4,1}{npnA}{nmos}{B}{C}{E}

\addResistor{mynand.out}{1.5,0}{b}{\orthogonalWireA}

\addPower{1,4}{power1}{$+5V$}
\addResistor{power1}{\Down}{npnA.B}{\orthogonalWireB}

\end{schema}
\end{subfigure}
\begin{subfigure}{.5\textwidth}
\centering
% \centering
% \begin{graphics}{0.9}{0.4}{-0.05}{80}{-1}{6}{t(ms)}{vs}{}
% \addPoints{blue}{(0,0)(10,0)(10,5)(15,5)(15,0)(20,0)(20,5)(25,5)(25,0)(30,0)(30,5)(35,5)(35,0)(100,0)}
% %\addLegend{Signal "KEY\_POWER"}
% \end{graphics}
\end{subfigure}

\begin{subfigure}{1\textwidth}
\centering
\begin{numeric}{Chronogramme du compteur 4 bits}
INPUT &  CC [blue]16{CC} CCC   \\
D0 &  HL 8{LHHL} LHL   \\
D1 &  H  4{LLLLHHHH} LLLL \\
D2 &  H 2{LLLLLLLLHHHHHHHH} LLLL   \\
D3 &  H{LLLLLLLLLLLLLLLLHHHHHHHHHHHHHHHH} LLLL  \\
END &  LL [green]14{LL} LHHLLL  \\
VALUE & L 2D{0} 2D{1} 2D{2} 2D{3} 2D{4} 2D{5} 2D{6} 2D{7} 2D{8} 2D{9} 2D{10} 2D{11} 2D{12} 2D{13} 2D{14} 2D{15} 2D{0} 2D{1}  \\
\end{numeric}%

\end{subfigure}
\caption{Titre}
\label{fig:fig}
\end{figure}


% \subsection{Affichage d'un graphique 2D avec insertion des données depuis plusieurs sources}


% \begin{graphics}{0.8}{0.4}{0}{20}{-1}{10}{x}{y}{Courbes de provenances diverses}
% \addPointsFromCSV{red}{comma}{src_examples/input_2.txt}
% \addTrace{green}{-10}{10}{x}
% \addPoints{blue}{(0,0)(10,0)(10,5)(15,5)(15,0)(20,0)(20,5)(25,5)(25,0)(30,0)(30,5)(35,5)(35,0)(100,0)}
% \addLegend{s1,s2,s3}
% \end{graphics}
% \index{addTrace}


% \section{Graphiques 3D}

% \subsection{Affichage d'un graphique 3D avec insertion des données depuis une équation}

% \plot{Titre 3D}{x*y*y}
% \index{plot}

% \chapter{Bilbiothèque 'Adding'}

% \section{Nomenclature}

% Un exemple de nomenclature est présenté à la page suivante
% \section{Exemples Layout}


% Une onomatopée\index{onomatopée} est une catégorie d'interjections\index{interjection}
% émise pour simuler un bruit particulier associé à un être,
% un animal ou un objet, par l'imitation des sons\index{son} que ceux-ci produisent

\nomenclature[E]{$r$}{Rapport cyclique d'un signal périodique}
\nomenclature[A]{$A_d$}{Coefficient d'amplification, gain différentiel }
\nomenclature[A]{$\varepsilon$}{Tension différentielle $(\varepsilon = E_+ - E_-)$\addUnit{V}}
\nomenclature[A]{$E_+$}{Tension entrée non inverseuse \addUnit{V}}
\nomenclature[A]{$E_-$}{Tension entrée inverseuse \addUnit{V}}
\nomenclature[E]{$\eta$}{Rendement d'un mécanisme \addUnit{\%}}
\nomenclature[E]{$\varphi$}{Déphasage entre deux signaux \addUnit{rad}}

\printnomenclature
\index{printnomenclature}
\index{nomenclature}
\index{addUnit}

\printindex 
\index{printindex}

\end{document}

\end{document}
%#############################################################
